\chapter{Ensiklopedia Master Bahasa Jepang (Final)}

% ================= THEME 1 =================
\begin{lessonbox}{1. Pengalaman Pertama \& Mengubah Kata Kerja Jadi Benda}
    \textbf{Topik: Using 初めて, Nominalization, Noun Reference (の), Qualifiers, Noun Suffixes (\#1-\#3, に)} \\
    Saat liburan ke luar negeri, banyak sekali "Pengalaman Pertama" yang akan kalian hadapi. Kata \textbf{初めて (Hajimete - Pertama kali)} adalah senjata rahasia turis!
    \begin{itemize}
        \item \textbf{Hajimete desu:} Kalau pelayan restoran bertanya apakah kamu sudah pernah ke sini, jawab saja \textit{"Hajimete desu"}. Mereka pasti akan menjelaskan cara pesannya dengan sangat ramah!
        \item \textbf{Nominalization (Menjadikan Benda):} Kalau kamu mau bilang "Makan sushi itu menyenangkan", kamu harus mengubah kata kerja "Makan" jadi benda dengan menambah \textbf{の (no)} atau \textbf{こと (koto)}. Contoh: \textit{Sushi o taberu \textbf{no} wa tanoshii desu!}
        \item \textbf{Noun Suffixes に (ni):} Partikel \textit{ni} juga sangat hebat. Bisa dipakai untuk tujuan: \textit{Kaimono \textbf{ni} ikimasu} (Pergi \textbf{untuk} belanja).
    \end{itemize}
\end{lessonbox}

\begin{convobox}{1}
    \noindent
    \jpword{にほん}{日本}{Nihon} \jpkana{は}{wa} 
    \jpword{はじ}{初}{haji}\jpkana{めてです。}{mete desu.}

    \vspace{1.5em}
    \noindent{\Large \color{articolor} \textbf{Arti:} Ini pertama kalinya (saya datang) ke Jepang.}
\end{convobox}

\begin{convobox}{2}
    \noindent
    \jpword{ふじさん}{富士山}{Fuji-san} \jpkana{に}{ni} 
    \jpword{のぼ}{登}{nobo}\jpkana{る}{ru} 
    \jpkana{の}{no} \jpkana{は}{wa} 
    \jpword{たの}{楽}{tano}\jpkana{しいです。}{shii desu.}

    \vspace{1.5em}
    \noindent{\Large \color{articolor} \textbf{Arti:} Mendaki Gunung Fuji ITU menyenangkan.}
\end{convobox}

% ================= THEME 2 =================
\begin{lessonbox}{2. Menjelaskan Alasan dengan Sangat Lancar}
    \textbf{Topik: Sentence Links (Because 1-3), わけ (Wake), Explanations, Giving Information} \\
    Kenapa keretanya terlambat? Kenapa makanannya enak? Orang Jepang sangat suka menjelaskan alasan!
    \begin{itemize}
        \item \textbf{〜から (〜kara) / 〜ので (〜node):} Ini adalah kata "Karena" yang paling standar. \textit{Ame ga futta node...} (Karena turun hujan...).
        \item \textbf{わけ (Wake):} Pernahkah makan ramen yang sangat enak sampai kamu berpikir "Pantas saja enak, ternyata kedai legendaris!"? Dalam bahasa Jepang, itu adalah \textbf{美味しいわけです (Oishii wake desu - Pantas saja enak! / Masuk akal kalau ini enak!)}. Ini akan membuat teman Jepangmu kagum!
    \end{itemize}
\end{lessonbox}

\begin{convobox}{3}
    \noindent
    \jpword{ゆき}{雪}{Yuki} \jpkana{が}{ga} 
    \jpword{ふ}{降}{fu}\jpkana{っている}{tte iru} 
    \jpword{わけ}{訳}{wake} \jpkana{ですね。}{desu ne.}

    \vspace{1.5em}
    \noindent{\Large \color{articolor} \textbf{Arti:} Pantas saja (masuk akal jika) turun salju ya. (Pantas saja dingin banget!)}
\end{convobox}

% ================= THEME 3 =================
\begin{lessonbox}{3. Memberi Penekanan, Contoh, dan Penyangkalan Total!}
    \textbf{Topik: Emphasis (こと/Other), Negation (~ない), も～ない, Negative Verb Endings, など・なんて・なんか, Giving Examples} \\
    Bagaimana cara bilang "Tidak ada SATU PUN"? Atau menyebutkan contoh oleh-oleh?
    \begin{itemize}
        \item \textbf{も〜ない (Mo...nai) / Penyangkalan Total:} Jika kamu ditanya apakah ada barang yang tertinggal, dan kamu yakin tidak ada, katakan \textbf{一つもないです (Hitotsu mo nai desu - Satu pun tidak ada)}.
        \item \textbf{など (Nado - Dan lain-lain):} Saat membeli banyak hal. \textit{Kashi ya, o-cha \textbf{nado} o kaimashita} (Saya membeli snack, teh, dan lain-lain).
        \item \textbf{なんか (Nanka):} Versi sangat kasual dari \textit{Nado}. Sering dipakai anak muda! \textit{Sushi nanka tabetai na~} (Pengen makan sushi atau semacamnya nih~).
    \end{itemize}
\end{lessonbox}

\begin{convobox}{4}
    \noindent
    \jpword{おかね}{お金}{O-kane} \jpkana{が}{ga} 
    \jpword{いちえん}{1円}{ichi-en} \jpkana{も}{mo} 
    \jpkana{ないです。}{nai desu.}

    \vspace{1.5em}
    \noindent{\Large \color{articolor} \textbf{Arti:} Tidak ada uang 1 yen pun (Sama sekali tidak punya uang).}
\end{convobox}

% ================= THEME 4 =================
\begin{lessonbox}{4. Kuantitas, Perbandingan, dan Derajat (Level)}
    \textbf{Topik: ほど expanded, Degree/Amount, Comparisons, (Not) Only} \\
    Adik-adik pasti suka membandingkan mainan atau makanan. Mari kita pelajari levelnya!
    \begin{itemize}
        \item \textbf{〜ほど〜ない (〜hodo...nai):} Artinya "Tidak se-...". Contoh: \textit{Oosaka wa, Toukyou \textbf{hodo} samuku nai desu} (Osaka, TIDAK SE-dingin Tokyo).
        \item \textbf{〜だけ (〜dake) / 〜しか〜ない (〜shika...nai):} \textit{Dake} artinya "Hanya". Tapi kalau mau lebih dramatis, pakai \textit{Shika + Tidak}. Contoh: \textit{Hyaku-en \textbf{shika nai} desu} (Hanya ada 100 yen saja / Tidak ada selain 100 yen).
    \end{itemize}
\end{lessonbox}

\begin{convobox}{5}
    \noindent
    \jpword{きょう}{今日}{Kyou} \jpkana{は}{wa} 
    \jpword{きのう}{昨日}{kinou} \jpkana{ほど}{hodo} 
    \jpword{あつ}{暑}{atsu}\jpkana{くないです。}{kunai desu.}

    \vspace{1.5em}
    \noindent{\Large \color{articolor} \textbf{Arti:} Hari ini tidak SE-panas kemarin.}
\end{convobox}

% ================= THEME 5 =================
\begin{lessonbox}{5. Permintaan Tingkat Lanjut \& Kebutuhan Mendesak}
    \textbf{Topik: Commands/Requests 1-4, Necessity} \\
    Kadang di bandara atau hotel, kamu harus mengisi formulir atau diminta melakukan sesuatu oleh petugas.
    \begin{itemize}
        \item \textbf{〜てください (〜te kudasai):} Tolong lakukan. (Bentuk standar).
        \item \textbf{〜ないでください (〜nai de kudasai):} Tolong JANGAN lakukan! \textit{Koko ni hairanai de kudasai!} (Tolong jangan masuk ke sini!).
        \item \textbf{Necessity (Harus):} Jika petugas bilang \textit{Kakanakereba narimasen}, artinya kamu "Harus / Wajib" menulis sesuatu di formulir tersebut!
    \end{itemize}
\end{lessonbox}

\begin{convobox}{6}
    \noindent
    \jpkana{ここで}{Koko de} 
    \jpword{た}{食}{ta}\jpkana{べないで}{benai de} 
    \jpword{くだ}{下}{kuda}\jpkana{さい。}{sai.}

    \vspace{1.5em}
    \noindent{\Large \color{articolor} \textbf{Arti:} Tolong JANGAN makan di sini.}
\end{convobox}

% ================= THEME 6 =================
\begin{lessonbox}{6. Harapan, Doa, dan Berandai-andai (If / Wishes)}
    \textbf{Topik: Wishes \& Hopes 1-2, Using ように 1-2, Using もし} \\
    Saat kamu berdoa di Kuil Jepang (Jinja), bagaimana cara meminta permohonan dalam bahasa Jepang?
    \begin{itemize}
        \item \textbf{〜ように (〜you ni):} Digunakan untuk berdoa! Gabungkan dengan tepuk tangan 2 kali di kuil dan ucapkan: \textit{Nihongo ga jouzu ni narimasu \textbf{you ni}!} (Semoga bahasa Jepangku jadi pintar!).
        \item \textbf{もし (Moshi):} Artinya "Seandainya/Misalkan". \textit{\textbf{Moshi} jikan ga attara, asobimashou} (Seandainya ada waktu, ayo kita main!).
    \end{itemize}
\end{lessonbox}

\begin{convobox}{7}
    \noindent
    \jpword{あした}{明日}{Ashita} 
    \jpword{は}{晴}{ha}\jpkana{れます}{remasu} 
    \jpkana{ように。}{you ni.}

    \vspace{1.5em}
    \noindent{\Large \color{articolor} \textbf{Arti:} Semoga besok cuacanya cerah.}
\end{convobox}

% ================= THEME 7 =================
% \begin{lessonbox}{7. Memanfaatkan Waktu: "Mumpung" \& Penyelesaian}
%     \textbf{Topik: Actions \& Time (うち, Endings, #3), Continuing Actions, Completing Actions 1-2} \\
%     Waktu liburan di Jepang sangat berharga! Mari kita manfaatkan "Mumpung"!
%     \begin{itemize}
%         \item \textbf{〜うちに (〜uchi ni):} Artinya "Selagi / Mumpung". \textit{Nihon ni iru \textbf{uchi ni}, takusan tabetai!} (Mumpung/Selagi masih ada di Jepang, aku mau makan banyak!).
%         \item \textbf{Completing Actions (〜てしまう / 〜てある):} \textit{Te shimau} (Sudah selesai semua / Tidak sengaja). Sedangkan \textit{Te aru} artinya suatu tindakan sudah disiapkan. \textit{Hoteru wa yoyaku shi\textbf{te arimasu}} (Hotelnya sudah dalam keadaan selesai dipesan).
%     \end{itemize}
% \end{lessonbox}

\begin{convobox}{8}
    \noindent
    \jpword{あたた}{温}{Atata}\jpkana{かい}{kai} 
    \jpword{うち}{内}{uchi} \jpkana{に}{ni} 
    \jpword{た}{食}{ta}\jpkana{べて}{bete} 
    \jpword{くだ}{下}{kuda}\jpkana{さい。}{sai.}

    \vspace{1.5em}
    \noindent{\Large \color{articolor} \textbf{Arti:} Tolong makan SELAGI (mumpung) masih hangat.}
\end{convobox}

% ================= THEME 8 =================
\begin{lessonbox}{8. Meskipun, Padahal, dan Sebagai Gantinya (Konjungsi Ajaib)}
    \textbf{Topik: Conjunctions でも 1-2, Although, Even Though, Instead, Considering} \\
    Ini adalah tata bahasa untuk anak pintar yang suka berdebat dengan alasan yang logis!
    \begin{itemize}
        \item \textbf{〜のに (〜noni):} Artinya "Padahal". Mengandung perasaan kecewa. \textit{Takai \textbf{noni}, oishikunai!} (PADAHAL harganya mahal, tapi tidak enak!).
        \item \textbf{〜かわりに (〜kawari ni):} Artinya "Sebagai gantinya / Alih-alih". \textit{Kyou wa sushi no \textbf{kawari ni}, ramen o tabemasu} (Hari ini sebagai ganti sushi, kita makan ramen).
    \end{itemize}
\end{lessonbox}

\begin{convobox}{9}
    \noindent
    \jpword{あめ}{雨}{Ame} \jpkana{が}{ga} 
    \jpword{ふ}{降}{fu}\jpkana{っている}{tte iru} 
    \jpkana{のに、}{noni,} 
    \jpkana{かさを}{kasa o} 
    \jpword{も}{持}{mo}\jpkana{っていません。}{tte imasen.}

    \vspace{1.5em}
    \noindent{\Large \color{articolor} \textbf{Arti:} PADAHAL sedang turun hujan, (tapi) saya tidak membawa payung.}
\end{convobox}

% ================= THEME 9 =================
\begin{lessonbox}{9. Melihat Sifat, Gosip, dan Kesan Orang Lain}
    \textbf{Topik: Using っぽい, Using ～がる, Seeming/Pretending, Reporting Information 1-2, Expectations} \\
    Menebak-nebak sifat barang atau perasaan orang Jepang lain di sekitarmu:
    \begin{itemize}
        \item \textbf{〜っぽい (〜ppoi):} Kelihatan seperti / Mirip seperti. \textit{Kodomo-ppoi} (Sifatnya seperti anak-anak).
        \item \textbf{〜がる (〜garu):} Kita tidak bisa menebak perasaan orang 100\%. Kalau ada temanmu yang menggigil, jangan bilang "Dia dingin", tapi bilangnya \textit{Samu\textbf{gatte iru}} (Dia KELIHATANNYA merasa dingin).
        \item \textbf{〜はずです (〜hazu desu):} Ekspektasi / Seharusnya. \textit{Densha wa 10-ji ni kuru \textbf{hazu desu}} (SEHARUSNYA keretanya datang jam 10).
    \end{itemize}
\end{lessonbox}

\begin{convobox}{10}
    \noindent
    \jpkana{あの}{Ano} \jpword{ひと}{人}{hito} \jpkana{は}{wa} 
    \jpword{さむ}{寒}{samu}\jpkana{がっています。}{gatte imasu.}

    \vspace{1.5em}
    \noindent{\Large \color{articolor} \textbf{Arti:} Orang itu (kelihatannya) sedang merasa kedinginan.}
\end{convobox}

\begin{convobox}{11}
    \noindent
    \jpword{あした}{明日}{Ashita} \jpkana{の}{no} 
    \jpword{てんき}{天気}{tenki} \jpkana{は}{wa} 
    \jpword{は}{晴}{ha}\jpkana{れる}{reru} 
    \jpword{はず}{筈}{hazu} \jpkana{です。}{desu.}

    \vspace{1.5em}
    \noindent{\Large \color{articolor} \textbf{Arti:} Cuaca besok SEHARUSNYA (pasti) cerah.}
\end{convobox}

% ================= THEME 10 =================
\begin{lessonbox}{10. Bahasa Gaul Tingkat Akhir (Informal Endings \& Particles Expanded)}
    \textbf{Topik: Informal Sentence Endings, More Questions, Particles Expanded} \\
    Tingkat tertinggi belajar bahasa Jepang adalah berbicara seperti karakter anime atau teman akrab yang sedang bersantai!
    \begin{itemize}
        \item \textbf{〜かな (〜kana):} Diucapkan saat berbicara sendiri atau ragu. \textit{Doko ni ikou kana~?} (Kira-kira kita pergi ke mana ya~?). Sangat imut untuk anak-anak!
        \item \textbf{〜さ (〜sa) / 〜ぜ (〜ze) / 〜わ (〜wa):} Partikel di akhir kalimat untuk gaya bicara tertentu. \textit{Ze} (Kasar/Laki-laki banget), \textit{Wa} (Feminin/Anggun), \textit{Sa} (Santai).
        \item \textbf{Expanded Particles:} Menggabungkan partikel! \textit{Tomodachi \textbf{to no} shashin} (Foto BERSAMA teman). \textit{Nihon \textbf{kara no} omiyage} (Oleh-oleh DARI Jepang).
    \end{itemize}
\end{lessonbox}

\begin{convobox}{12}
    \noindent
    \jpword{きょう}{今日}{Kyou} \jpkana{は}{wa} 
    \jpword{なに}{何}{nani} \jpkana{を}{o} 
    \jpword{た}{食}{ta}\jpkana{べよう}{beyou} \jpkana{かな〜。}{kana~.}

    \vspace{1.5em}
    \noindent{\Large \color{articolor} \textbf{Arti:} Hari ini, kira-kira enaknya makan apa ya~.}
\end{convobox}

% --- SUMMARY BOX CLOSING ---
\vspace{2em}
\begin{tcolorbox}[
    colback=yellow!5!white,
    colframe=orange!80!black,
    boxrule=4pt,
    arc=6mm,
    title={\huge\bfseries \sffamily 🎆 T A M A T - KELULUSAN MASTER BAHASA JEPANG! 🎆},
    fonttitle=\color{white},
    attach boxed title to top center={yshift=-5mm},
    boxed title style={colback=orange!80!black, rounded corners},
    left=5mm, right=5mm, top=7mm, bottom=5mm
]
\Large \textbf{Air Mata Kebahagiaan! Kalian Berhasil Menyelesaikannya!}\\[0.5em]
Ayah, Bunda, dan Adik-adik yang super hebat, dengan selesainya Bab Puncak (Revisit 4) ini, \textbf{seluruh materi di dalam Andi - Japanese Handbook telah tamat 100\%!} 

Dari tidak bisa mengucapkan sepatah kata pun, hingga bisa memesan Ramen di mesin tiket, mengobrol dengan kasir Konbini, menebak perasaan orang dengan \textit{\textasciitilde garu}, hingga mendoakan cuaca di kuil dengan \textit{\textasciitilde you ni}. Ini adalah pencapaian yang luar biasa!

\textbf{Pesan Terakhir untuk Turis Hebat:}
\begin{itemize}
    \item Jangan takut berbuat salah saat berbicara di Jepang. Orang Jepang sangat menghargai senyuman dan usahamu!
    \item Bawalah buku ini kemanapun kalian pergi di Jepang. Jadikan buku ini saksi petualangan keluarga kalian yang luar biasa.
\end{itemize}

\begin{center}
    \Huge \textbf{いってらっしゃい！}\\[0.2em]
    (Itterasshai!)\\[0.2em]
    \Large Selamat Berpetualang di Jepang!
\end{center}
\end{tcolorbox}

\newpage